\section{Integration by Parts and Boundary Terms}
\label{sec:integration_by_parts_boundary_terms}

Integration by parts is the operation that converts the variation of a derivative into the derivative of a coefficient. In variational mechanics, this step separates the first variation into two distinct pieces:

\begin{itemize}
    \item an interior term that produces the differential equation of motion;
    \item a boundary term that records what happens at the endpoints.
\end{itemize}

The same structure reappears in every later field-theoretic derivation. Understanding it in one time dimension is therefore more important than memorizing the final Euler--Lagrange formula.

\subsection{The ordinary integration-by-parts identity}

For sufficiently differentiable functions \(f(t)\) and \(g(t)\),

\[
\int_{t_{1}}^{t_{2}}
f(t)\dot g(t)
\,\mathrm{d}t
=
\left[f(t)g(t)\right]_{t_{1}}^{t_{2}}
-
\int_{t_{1}}^{t_{2}}
\dot f(t)g(t)
\,\mathrm{d}t.
\]

This follows from the product rule

\[
\frac{\mathrm{d}}{\mathrm{d}t}(fg)
=
\dot f g+f\dot g.
\]

Rearranging,

\[
f\dot g
=
\frac{\mathrm{d}}{\mathrm{d}t}(fg)-\dot f g,
\]

and integrating over the interval gives the identity.

The notation

\[
[fg]_{t_{1}}^{t_{2}}
=
f(t_{2})g(t_{2})-f(t_{1})g(t_{1})
\]

makes the endpoint contribution explicit.

\subsection{Application to the first variation}

For the action

\[
S[q]
=
\int_{t_{1}}^{t_{2}}
L(q,\dot q,t)
\,\mathrm{d}t,
\]

the first variation is

\[
\delta S
=
\int_{t_{1}}^{t_{2}}
\left(
\frac{\partial L}{\partial q}\delta q
+
\frac{\partial L}{\partial\dot q}\delta\dot q
\right)
\,\mathrm{d}t.
\]

Since

\[
\delta\dot q
=
\frac{\mathrm{d}}{\mathrm{d}t}(\delta q),
\]

the second term has the form required for integration by parts. Set

\[
f(t)
=
\frac{\partial L}{\partial\dot q},
\qquad
g(t)=\delta q(t).
\]

Then

\[
\int_{t_{1}}^{t_{2}}
\frac{\partial L}{\partial\dot q}
\delta\dot q
\,\mathrm{d}t
=
\left[
\frac{\partial L}{\partial\dot q}
\delta q
\right]_{t_{1}}^{t_{2}}
-
\int_{t_{1}}^{t_{2}}
\frac{\mathrm{d}}{\mathrm{d}t}
\left(
\frac{\partial L}{\partial\dot q}
\right)
\delta q
\,\mathrm{d}t.
\]

Therefore

\[
\boxed{
\delta S
=
\left[
\frac{\partial L}{\partial\dot q}
\delta q
\right]_{t_{1}}^{t_{2}}
+
\int_{t_{1}}^{t_{2}}
\left[
\frac{\partial L}{\partial q}
-
\frac{\mathrm{d}}{\mathrm{d}t}
\left(
\frac{\partial L}{\partial\dot q}
\right)
\right]
\delta q
\,\mathrm{d}t
}.
\]

The derivative has been transferred from \(\delta q\) to \(\partial L/\partial\dot q\). This is the decisive step in the derivation.

\subsection{Why the derivative must be moved}

Before integration by parts, the first variation contains both \(\delta q\) and \(\delta\dot q\). These are not independent functions because

\[
\delta\dot q
=
\frac{\mathrm{d}}{\mathrm{d}t}(\delta q).
\]

The fundamental lemma of the calculus of variations applies to an integral of the form

\[
\int f(t)\delta q(t)\,\mathrm{d}t.
\]

Integration by parts rewrites the complete interior contribution in precisely this form. The arbitrariness of \(\delta q\) can then be used to conclude that its coefficient vanishes.

Without this step, one cannot correctly infer separate conditions from the coefficients of \(\delta q\) and \(\delta\dot q\).

\subsection{The boundary term as a separate object}

The boundary contribution is

\[
B
=
\left[
\frac{\partial L}{\partial\dot q}
\delta q
\right]_{t_{1}}^{t_{2}}.
\]

Using the canonical momentum

\[
p
=
\frac{\partial L}{\partial\dot q},
\]

we may write

\[
B=[p\,\delta q]_{t_{1}}^{t_{2}}.
\]

This term contains no time integral. It depends only on endpoint data.

For fixed endpoint positions,

\[
\delta q(t_{1})=
\delta q(t_{2})=0,
\]

so

\[
B=0.
\]

For free endpoints, the term survives and produces natural boundary conditions. Thus one should never write ``the boundary term is zero'' without first stating the conditions that make it zero.

\subsection{A direct example}

Consider

\[
L
=
\frac{m}{2}\dot q^{2}-V(q).
\]

The first variation is

\[
\delta S
=
\int_{t_{1}}^{t_{2}}
\left(
-V'(q)\delta q
+
m\dot q\,\delta\dot q
\right)
\,\mathrm{d}t.
\]

Integrating the second term by parts,

\[
\int_{t_{1}}^{t_{2}}
m\dot q\,\delta\dot q
\,\mathrm{d}t
=
[m\dot q\,\delta q]_{t_{1}}^{t_{2}}
-
\int_{t_{1}}^{t_{2}}
m\ddot q\,\delta q
\,\mathrm{d}t.
\]

Hence

\[
\delta S
=
[m\dot q\,\delta q]_{t_{1}}^{t_{2}}
-
\int_{t_{1}}^{t_{2}}
\left(
m\ddot q+V'(q)
\right)
\delta q
\,\mathrm{d}t.
\]

For fixed endpoints, the first term vanishes. Stationarity then implies

\[
m\ddot q+V'(q)=0.
\]

Every sign in this equation can be traced back to the product rule and the minus sign introduced by integration by parts.

\subsection{Sign checking}

A reliable way to check the sign is to start from

\[
\frac{\mathrm{d}}{\mathrm{d}t}
\left(
\frac{\partial L}{\partial\dot q}
\delta q
\right)
=
\frac{\mathrm{d}}{\mathrm{d}t}
\left(
\frac{\partial L}{\partial\dot q}
\right)
\delta q
+
\frac{\partial L}{\partial\dot q}
\delta\dot q.
\]

Solving for the last term,

\[
\frac{\partial L}{\partial\dot q}
\delta\dot q
=
\frac{\mathrm{d}}{\mathrm{d}t}
\left(
\frac{\partial L}{\partial\dot q}
\delta q
\right)
-
\frac{\mathrm{d}}{\mathrm{d}t}
\left(
\frac{\partial L}{\partial\dot q}
\right)
\delta q.
\]

The minus sign is therefore unavoidable.

\subsection{Adding a total time derivative to the Lagrangian}

Suppose two Lagrangians differ by a total derivative:

\[
L'(q,\dot q,t)
=
L(q,\dot q,t)
+
\frac{\mathrm{d}}{\mathrm{d}t}F(q,t).
\]

Their actions satisfy

\[
\begin{aligned}
S'[q]
&=
\int_{t_{1}}^{t_{2}}
L'\,\mathrm{d}t\\
&=
S[q]
+
\int_{t_{1}}^{t_{2}}
\frac{\mathrm{d}F}{\mathrm{d}t}
\,\mathrm{d}t\\
&=
S[q]
+
F(q(t_{2}),t_{2})
-
F(q(t_{1}),t_{1}).
\end{aligned}
\]

Thus the two actions differ only by an endpoint term.

For fixed endpoint values and fixed endpoint times, the variation of this additional term vanishes. Therefore \(L\) and \(L'\) produce the same Euler--Lagrange equations.

This establishes an important equivalence:

\[
\boxed{
L
\sim
L+
\frac{\mathrm{d}F}{\mathrm{d}t}
}.
\]

The Lagrangian is not unique. Different Lagrangians can encode the same equations of motion when they differ by a total derivative.

\subsection{Example of equivalent Lagrangians}

Consider

\[
L
=
\frac{m}{2}\dot q^{2}-V(q)
\]

and choose

\[
F(q,t)=\alpha q t,
\]

where \(\alpha\) is constant. Then

\[
\frac{\mathrm{d}F}{\mathrm{d}t}
=
\alpha q+\alpha t\dot q.
\]

The modified Lagrangian is

\[
L'
=
\frac{m}{2}\dot q^{2}-V(q)
+
\alpha q+\alpha t\dot q.
\]

Although the individual partial derivatives of \(L'\) differ from those of \(L\), the extra contributions cancel in the Euler--Lagrange equation.

Indeed,

\[
\frac{\partial L'}{\partial q}
=
-V'(q)+\alpha,
\]

and

\[
\frac{\partial L'}{\partial\dot q}
=
m\dot q+\alpha t.
\]

Therefore

\[
\frac{\mathrm{d}}{\mathrm{d}t}
\left(
\frac{\partial L'}{\partial\dot q}
\right)
=
m\ddot q+\alpha.
\]

The Euler--Lagrange equation becomes

\[
m\ddot q+\alpha-(-V'(q)+\alpha)=0,
\]

so

\[
m\ddot q+V'(q)=0.
\]

The dynamics is unchanged.

\subsection{Boundary terms are not always negligible}

A boundary term may be physically or mathematically important when

\begin{itemize}
    \item endpoint values are not fixed;
    \item the integration region has a physical boundary;
    \item one studies conserved charges or canonical momenta;
    \item the action contains higher derivatives;
    \item the theory is defined on a noncompact region and fall-off conditions matter;
    \item surface contributions encode measurable quantities.
\end{itemize}

Thus ``boundary term'' does not mean ``unimportant term.'' It means a term supported on the boundary rather than in the interior.

\subsection{Repeated integration by parts}

If an action depends on higher derivatives, for example

\[
S[q]
=
\int
L(q,\dot q,\ddot q,t)
\,\mathrm{d}t,
\]

then varying \(\ddot q\) produces

\[
\delta\ddot q
=
\frac{\mathrm{d}^{2}}{\mathrm{d}t^{2}}(\delta q).
\]

Two integrations by parts are required to remove derivatives from \(\delta q\). The resulting boundary terms contain both \(\delta q\) and \(\delta\dot q\).

This explains why higher-derivative actions generally require more boundary data and produce higher-order equations of motion.

The present course will mainly use first-derivative Lagrangians, but the pattern is worth recognizing.

\subsection{Several coordinates}

For generalized coordinates \(q^{i}\),

\[
\delta S
=
\int_{t_{1}}^{t_{2}}
\sum_{i}
\left(
\frac{\partial L}{\partial q^{i}}\delta q^{i}
+
\frac{\partial L}{\partial\dot q^{i}}\delta\dot q^{i}
\right)
\,\mathrm{d}t.
\]

Integrating each velocity term by parts gives

\[
\delta S
=
\left[
\sum_{i}
\frac{\partial L}{\partial\dot q^{i}}
\delta q^{i}
\right]_{t_{1}}^{t_{2}}
+
\int_{t_{1}}^{t_{2}}
\sum_{i}
\left[
\frac{\partial L}{\partial q^{i}}
-
\frac{\mathrm{d}}{\mathrm{d}t}
\left(
\frac{\partial L}{\partial\dot q^{i}}
\right)
\right]
\delta q^{i}
\,\mathrm{d}t.
\]

The boundary term is a sum of momentum-coordinate pairings:

\[
\left[
\sum_{i}p_{i}\delta q^{i}
\right]_{t_{1}}^{t_{2}}.
\]

\subsection{The field-theory analogue}

For a scalar field with action

\[
S[\phi]
=
\int_{\Omega}
\mathcal{L}(\phi,\partial_{\mu}\phi,x)
\,\mathrm{d}^{4}x,
\]

the first variation contains

\[
\int_{\Omega}
\frac{\partial\mathcal{L}}
{\partial(\partial_{\mu}\phi)}
\partial_{\mu}(\delta\phi)
\,\mathrm{d}^{4}x.
\]

Spacetime integration by parts gives

\[
\begin{aligned}
&\int_{\Omega}
\frac{\partial\mathcal{L}}
{\partial(\partial_{\mu}\phi)}
\partial_{\mu}(\delta\phi)
\,\mathrm{d}^{4}x\\
&\qquad=
\int_{\partial\Omega}
\frac{\partial\mathcal{L}}
{\partial(\partial_{\mu}\phi)}
\delta\phi\,n_{\mu}
\,\mathrm{d}\Sigma
-
\int_{\Omega}
\partial_{\mu}
\left(
\frac{\partial\mathcal{L}}
{\partial(\partial_{\mu}\phi)}
\right)
\delta\phi
\,\mathrm{d}^{4}x.
\end{aligned}
\]

The endpoint term of mechanics becomes a surface integral over the boundary \(\partial\Omega\). The one-dimensional formula is therefore not merely analogous to the field calculation; it is its simplest special case.

\subsection{A diagnostic procedure}

Whenever integration by parts appears in a variational derivation:

\begin{enumerate}
    \item identify the derivative acting on the variation;
    \item write the product-rule identity explicitly;
    \item separate the boundary and interior terms;
    \item state the boundary conditions;
    \item justify whether the boundary term vanishes;
    \item check the sign of the interior derivative;
    \item only then apply the fundamental lemma.
\end{enumerate}

Following this sequence prevents hidden assumptions and sign errors.

\subsection{What has been established}

Integration by parts converts the derivative of a variation into an interior derivative acting on the canonical momentum and an explicit boundary contribution. The interior term yields the Euler--Lagrange equation, while the boundary term records the endpoint assumptions.

Lagrangians that differ by a total time derivative produce actions that differ only by endpoint terms and therefore give the same fixed-endpoint equations of motion.

The next section applies the full method to several concrete particle models, from free motion to oscillation and constrained dynamics.